% !TeX root = part6.tex
\documentclass[../final.tex]{subfiles}

\begin{document}

% Конспект: Desktop/1.jpg--3.jpg (10.10). В тетради a_k, \hbar=1;
% здесь обозначения part5 (b_k) и явная \hbar.
% Вариант по плану семинаров: LaTeX/archive_plan_2026-11-01.
\practice{6}{10.10.2026}{Фононы как идеальный бозе-газ}

\rtheory{Продолжение задачи о квантовании цепочки}

Продолжаем задачу~4.18 из предыдущего занятия. Напомним обозначения:
$N$ атомов массы $m$, жёсткость пружин $\alpha$, безразмерное волновое
число $k=2\pi j/N$ и
\[
    \hat Q_k=\frac1{\sqrt N}\sum_r\hat q_re^{-ikr},\qquad
    \hat P_k=\frac1{\sqrt N}\sum_r\hat p_re^{ikr},\qquad
    [\hat q_r,\hat p_s]=i\hbar\delta_{rs}.
\]
Гамильтониан в нормальных координатах:
\[
    \hat H=\frac1{2m}\sum_k\hat P_k\hat P_{-k}
       +\alpha\sum_k(1-\cos k)\hat Q_k\hat Q_{-k},
    \qquad
    \alpha(1-\cos k)=\frac{m\omega_k^2}{2},
\]
то есть $\omega_k^2=(4\alpha/m)\sin^2(k/2)$.
Операторы уничтожения и рождения фонона ($k\ne0$):
\begin{align*}
    \hat b_k&=\frac1{\sqrt{2\hbar m\omega_k}}
       \left(m\omega_k\hat Q_k+i\hat P_{-k}\right),\\
    \hat b_k^\dagger&=\frac1{\sqrt{2\hbar m\omega_k}}
       \left(m\omega_k\hat Q_{-k}-i\hat P_k\right).
\end{align*}

\subsection*{Проверка коммутатора}

Из $[\hat q_r,\hat p_s]=i\hbar\delta_{rs}$ следует
$[\hat Q_k,\hat P_{k'}]=i\hbar\delta_{kk'}$, а координаты и импульсы
между собой коммутируют. Поэтому в $[\hat b_k,\hat b_{k'}^\dagger]$
остаются только перекрёстные члены:
\begin{align*}
    [\hat b_k,\hat b_{k'}^\dagger]
       &=\frac1{2\hbar m\sqrt{\omega_k\omega_{k'}}}
         \Bigl(-im\omega_k[\hat Q_k,\hat P_{k'}]
               +im\omega_{k'}[\hat P_{-k},\hat Q_{-k'}]\Bigr)\\
       &=\frac{\hbar m(\omega_k+\omega_{k'})}
              {2\hbar m\sqrt{\omega_k\omega_{k'}}}\,\delta_{kk'}
        =\delta_{kk'}.
\end{align*}
Аналогично $[\hat b_k,\hat b_{k'}]=0$: два слагаемых сокращаются,
поскольку $\omega_k=\omega_{-k}$.

\subsection*{Гамильтониан в операторах рождения и уничтожения}

Обратные формулы:
\[
    \hat Q_k=\sqrt{\frac{\hbar}{2m\omega_k}}
       \left(\hat b_k+\hat b_{-k}^\dagger\right),\qquad
    \hat P_k=\frac{\sqrt{2\hbar m\omega_k}}{2i}
       \left(\hat b_{-k}-\hat b_k^\dagger\right).
\]
Подставляя их в $\hat H$, получаем
\begin{align*}
    \hat H&=\sum_k\frac{\hbar\omega_k}4\Bigl[
       -\left(\hat b_{-k}-\hat b_k^\dagger\right)
        \left(\hat b_k-\hat b_{-k}^\dagger\right)
       +\left(\hat b_k+\hat b_{-k}^\dagger\right)
        \left(\hat b_{-k}+\hat b_k^\dagger\right)\Bigr]\\
    &=\sum_k\frac{\hbar\omega_k}4\Bigl[
       \hat b_k\hat b_k^\dagger+\hat b_k^\dagger\hat b_k
       +\hat b_{-k}\hat b_{-k}^\dagger+\hat b_{-k}^\dagger\hat b_{-k}\Bigr].
\end{align*}
Члены $\hat b_k\hat b_{-k}$ и $\hat b_k^\dagger\hat b_{-k}^\dagger$
сократились, так как эти операторы коммутируют.
После замены $-k\to k$ во второй паре и использования
$\hat b_k\hat b_k^\dagger=\hat b_k^\dagger\hat b_k+1$:
\[
    \hat H=\sum_k\frac{\hbar\omega_k}4
       \left(2\hat b_k^\dagger\hat b_k+2+2\hat b_k^\dagger\hat b_k\right)
     =\sum_k\hbar\omega_k\left(\hat b_k^\dagger\hat b_k+\frac12\right).
\]
Собственные состояния --- $|\ldots n_j\ldots\rangle$, где $n_j$ ---
число фононов в моде $j$, и
\[
    \hat b_j|n_j\rangle=\sqrt{n_j}\,|n_j-1\rangle,\qquad
    \hat b_j^\dagger|n_j\rangle=\sqrt{n_j+1}\,|n_j+1\rangle.
\]

\rtheory{Тепловые средние}

Равновесное среднее оператора $\hat L$ вычисляется с распределением
Гиббса в базисе чисел заполнения:
\[
    \langle\hat L\rangle=
    \frac{\sum_{\{n\}}\langle\ldots n_j\ldots|\hat L\,
          e^{-\hat H/T}|\ldots n_j\ldots\rangle}
         {\sum_{\{n\}}\langle\ldots n_j\ldots|
          e^{-\hat H/T}|\ldots n_j\ldots\rangle},
    \qquad
    \hat H=\sum_i\hbar\omega_i\left(\hat n_i+\frac12\right).
\]
Полагаем \rassumption{$k_{\mathrm B}=1$}.

\subsection*{Недиагональные средние}

Операторы $\hat b_j$ и $\hat b_j^\dagger$ меняют $n_j$ на единицу,
а след берётся по диагональным элементам. Поэтому
\[
    \langle\hat b_j\rangle=\langle\hat b_j^\dagger\rangle=0,\qquad
    \langle\hat b_j\hat b_{j'}\rangle
      =\langle\hat b_j^\dagger\hat b_{j'}^\dagger\rangle=0,\qquad
    \langle\hat b_j^\dagger\hat b_{j'}\rangle=0\quad(j\ne j').
\]

\subsection*{Число заполнения}

Ненулевым остаётся только $\langle\hat b_j^\dagger\hat b_j\rangle=\langle\hat n_j\rangle$.
Множители от остальных мод и от нулевой энергии сокращаются.
Обозначим $x=e^{-\hbar\omega_j/T}$:
\[
    \langle\hat n_j\rangle
       =\frac{\sum_{n=0}^\infty n\,e^{-\hbar\omega_j(n+1/2)/T}}
             {\sum_{n=0}^\infty e^{-\hbar\omega_j(n+1/2)/T}}
       =\frac{x/(1-x)^2}{1/(1-x)}
       =\frac1{e^{\hbar\omega_j/T}-1}.
\]
Это \rconcept{распределение Бозе--Эйнштейна} с нулевым химическим
потенциалом: \rresult{фононы --- бозоны, число которых не сохраняется}.
Средняя энергия решётки
\[
    E=\sum_k\hbar\omega_k\left(\bar n_k+\frac12\right)
     =\sum_k\frac{\hbar\omega_k}2\,
        \frac{e^{\hbar\omega_k/T}+1}{e^{\hbar\omega_k/T}-1}
     =\sum_k\frac{\hbar\omega_k}2\coth\frac{\hbar\omega_k}{2T}.
\]

\rtheory{То же распределение из подсчёта состояний}

Получим $\bar n$ микроканонически. Рассмотрим $N\gg1$ осцилляторов
одной частоты $\omega$, между которыми распределено $L\gg1$ квантов:
\[
    E=\hbar\omega\left(L+\frac N2\right).
\]

\begin{rbox}[left=18pt,right=18pt,top=12pt,bottom=10pt]{[]}
\centering
\begin{tikzpicture}[>=Latex,every node/.style={text=rtext}]
    % Последовательность: oo|ooo||o|o| --- 7 квантов, 5 перегородок.
    \foreach \i in {0,1,3,4,5,8,10}
        \fill[rc3] (0.5*\i,0) circle (0.14);
    \foreach \i in {2,6,7,9,11}
        \draw[rconceptcolor,line width=1.2pt] (0.5*\i,-0.35)--(0.5*\i,0.35);
    \draw[<->] (-0.2,-0.6)--(5.7,-0.6)
        node[midway,below] {$L+N-1$ мест};
\end{tikzpicture}
\captionof{figure}{Кванты (кружки) разделены перегородками на $N$
осцилляторов. Здесь $N=6$, $L=7$; между соседними перегородками
может не быть ни одного кванта.}
\end{rbox}

Каждое распределение задаётся расстановкой $L$ одинаковых квантов
и $N-1$ одинаковых перегородок, поэтому статистический вес
\[
    \Gamma=\frac{(L+N-1)!}{L!\,(N-1)!},\qquad S=\ln\Gamma.
\]
По формуле Стирлинга $\ln n!\simeq n\ln(n/e)$:
\[
    S\simeq(L+N-1)\ln\frac{L+N-1}{e}-L\ln\frac Le
       -(N-1)\ln\frac{N-1}{e}.
\]
Температура определяется из $dE=T\,dS$:
\[
    \frac1T=\frac{dS}{dE}=\frac1{\hbar\omega}\frac{dS}{dL}
       =\frac1{\hbar\omega}\ln\frac{L+N}{L}.
\]
Отсюда
\[
    \frac LN=\frac1{e^{\hbar\omega/T}-1},\qquad
    E=N\hbar\omega\left(\frac12+\frac1{e^{\hbar\omega/T}-1}\right).
\]
\rresult{Подсчёт состояний даёт то же планковское среднее число квантов
на осциллятор}, что и усреднение с распределением Гиббса.

\end{document}
